\subsection*{Layer 6b --- Yukawa Stage B (algebraic content)}

\begin{theorem}[\texttt{YukawaTopForm.finrank\_lambda3\_eq\_one}]
A \texttt{YukawaTopForm K E}, defined as a linear equivalence
$\Lambda^3 E \simeq K$, forces $\dim_K (\Lambda^3 E) = 1$.
\end{theorem}

\begin{theorem}[\texttt{YukawaTopForm.finrank\_E\_eq\_three}]
For a finite-dimensional $K$-vector space, a \texttt{YukawaTopForm}
forces $\dim_K E = 3$. The proof is the composition
\[
\Lambda^3 E \simeq K
\;\Rightarrow\; \dim \Lambda^3 E = 1
\;\xRightarrow{\text{\texttt{finrank\_eq}}}\; \binom{\dim E}{3} = 1
\;\xRightarrow{\text{\texttt{choose\_eq\_one\_iff}}}\; \dim E = 3.
\]
\end{theorem}

\begin{theorem}[\texttt{YukawaTopForm.ofFinrankEqThree}]
Conversely, any finite-dimensional $K$-vector space with $\dim E = 3$
admits a \texttt{YukawaTopForm}: pick any basis and read off the
resulting one-dimensional $\Lambda^3 E$ as $K$.
\end{theorem}

\begin{theorem}[\texttt{YukawaTopForm.nonempty\_iff\_finrank\_eq\_three}]
Characterisation: for a finite-dimensional $K$-vector space,
\[
\text{Nonempty}\,(\texttt{YukawaTopForm}\,K\,E)
\;\iff\;
\dim_K E = 3.
\]
The Stage B hypothesis is therefore \emph{precisely} the rank-3
condition, expressed in exterior-algebra form.
\end{theorem}

\begin{theorem}[\texttt{YukawaTopForm.toYukawaRankCertificate}]
A Stage B \texttt{YukawaTopForm} produces a Stage A
\texttt{YukawaRankCertificate} by extracting $\dim E = 3$ and
choosing an arbitrary basis.
\end{theorem}

\subsection*{Layer 6c --- Yukawa Stage C (primitive indecomposable
coupling)}

\begin{theorem}[\texttt{PrimitiveIndecomposableYukawaCoupling.finrank\_E\_eq\_three}]
A \texttt{PrimitiveIndecomposableYukawaCoupling K E}, defined as
a linear equivalence $E \simeq (\Lambda^2 E)^*$ plus
non-triviality of $E$, forces $\dim_K E = 3$.

\smallskip\noindent
\emph{Proof sketch.}\ The contraction iso transports dimensions
as
\[
\dim E
  = \dim (\Lambda^2 E)^*
  \overset{\text{dual}}{=} \dim \Lambda^2 E
  \overset{\text{exteriorPower}}{=} \binom{\dim E}{2}
  = \tfrac{n(n-1)}{2}.
\]
The integer equation $n = n(n-1)/2$ with $n \ge 1$ has the
unique solution $n = 3$ (a finite case split / \texttt{omega}).
\end{theorem}

\begin{theorem}[\texttt{PrimitiveIndecomposableYukawaCoupling.toYukawaTopForm}]
A Stage C coupling produces a Stage B
\texttt{YukawaTopForm} via the v4 converse
\texttt{YukawaTopForm.ofFinrankEqThree} applied to
\texttt{finrank\_E\_eq\_three}.
\end{theorem}

\begin{theorem}[\texttt{PrimitiveIndecomposableYukawaCoupling.toYukawaRankCertificate}]
A Stage C coupling produces a Stage A \texttt{YukawaRankCertificate}
through the composition C $\Rightarrow$ B $\Rightarrow$ A.
\end{theorem}

\begin{theorem}[\texttt{PrimitiveIndecomposableYukawaCoupling.ofFinrankEqThree}]
Conversely, any finite-dimensional $K$-vector space with
$\dim E = 3$ admits a Stage C coupling: a contraction iso
$E \simeq (\Lambda^2 E)^*$ exists by dimension count, since both
sides have dimension $3$ (via
\texttt{dual\_finrank\_eq} and \texttt{exteriorPower.finrank\_eq}).
\end{theorem}

\begin{theorem}[\texttt{PrimitiveIndecomposableYukawaCoupling.nonempty\_iff\_finrank\_eq\_three}]
For a finite-dimensional $K$-vector space,
$\text{Nonempty}\,(\texttt{PrimitiveIndecomposableYukawaCoupling}\,K\,E) \iff \dim_K E = 3$.
\end{theorem}

\subsection*{Layer 6d --- Yukawa Stage D (primary trilinear form)}

\begin{theorem}[\texttt{YukawaTrilinearForm.contraction}]
For every alternating trilinear form
$\omega : \Lambda^3 E \to K$, the associated contraction map
\[
C_\omega : E \to (\Lambda^2 E)^*,
\qquad
C_\omega(v)(w \wedge x) := \omega(v \wedge w \wedge x),
\]
is constructible as a \texttt{LinearMap}. The Lean construction
uses \texttt{exteriorPower.alternatingMapLinearEquiv} (to
convert between $(\Lambda^n E)^*$ and the space of $K$-valued
alternating maps on $E^n$) and \texttt{AlternatingMap.curryLeft}
(to split off the first argument).
\end{theorem}

\begin{theorem}[\texttt{YukawaTrilinearForm.ContractionInjective} and \texttt{ContractionSurjective}]
Two explicit operational \texttt{Prop} predicates on $\omega$:
\[
\texttt{ContractionInjective}\,\omega \;:=\; \texttt{Injective}\,C_\omega,
\qquad
\texttt{ContractionSurjective}\,\omega \;:=\; \texttt{Surjective}\,C_\omega.
\]
The TFPT-side phrases "non-degenerate" and "primitive / no
spectator" are docstring synonyms for the two operational
predicates respectively. The combined predicate
\texttt{PrimitiveIndecomposableYukawaCondition} bundles both
(plus \texttt{Nontrivial}\,$E$ at the use site).
\end{theorem}

\begin{theorem}[\texttt{YukawaTrilinearForm.contractionEquivOfBoth}]
A trilinear form $\omega$ satisfying both predicates over a
finite-dimensional vector space induces a linear iso
$E \simeq (\Lambda^2 E)^*$ --- the contraction $C_\omega$ itself,
upgraded to a \texttt{LinearEquiv} via
\texttt{LinearEquiv.ofBijective}.
\end{theorem}

\begin{theorem}[\texttt{YukawaTrilinearForm.toPrimitiveIndecomposableYukawaCoupling}]
A trilinear form with \texttt{ContractionInjective} and
\texttt{ContractionSurjective} on a non-trivial finite-dimensional
vector space yields a Stage C
\texttt{PrimitiveIndecomposableYukawaCoupling}. Together with the
Stage C $\Rightarrow$ Stage B $\Rightarrow$ Stage A chain, this means

\[
\boxed{\;
\omega \,:\,\Lambda^3 E \to K
\;\;\text{(alt.)}\;\;\text{plus}\;\;\texttt{CI}, \texttt{CS}, \texttt{Nontrivial}\,E
\;\;\Longrightarrow\;\;
\dim E = 3.
\;}
\]
\end{theorem}

\begin{theorem}[\texttt{YukawaTrilinearForm.finrank\_E\_eq\_zero\_or\_three} (v9)]
\emph{Unconditional transparency form, no \texttt{Nontrivial}
hypothesis.} For any finite-dimensional $K$-vector space $E$,
a trilinear form $\omega$ satisfying \texttt{ContractionInjective}
and \texttt{ContractionSurjective} forces
\[
\dim_K E = 0 \;\;\lor\;\; \dim_K E = 3.
\]
The two predicates are vacuously satisfied when $E = 0$;
the rank-3 conclusion holds precisely when $E$ is non-trivial.
This theorem records the fact that the operational predicates
do \emph{not} themselves rule out the trivial case, which is the
mathematical reason the rank-3 slogan above carries the
\texttt{Nontrivial}\,$E$ hypothesis.
\end{theorem}

\textbf{What Stage D says about TFPT Paper~2.}
The phrase "primitive indecomposable Yukawa type" in TFPT
Paper~2 contains three semi-formal qualifiers:

\begin{itemize}[leftmargin=1.6em, itemsep=0pt]
\item \emph{primitive} --- the trilinear coupling does not
    factor through a sub-trilinear form;
\item \emph{indecomposable} --- the carrier does not split as
    a direct sum that respects the coupling;
\item \emph{no spectator} --- the closure
    $E_- \otimes \Lambda^2 E_- \to \Lambda^3 E_-$ exhausts
    $E_-$.
\end{itemize}

\noindent
\emph{Lean translation in Stage D.}\ The primary datum is the
alternating trilinear form $\omega$ on $E_-$. The three
qualifiers translate to two explicit operational \texttt{Prop}
predicates on $\omega$:
\texttt{ContractionInjective} (the contraction is injective;
equivalently, no nonzero $v$ has $C_\omega(v) = 0$ ---
\emph{non-degenerate} in TFPT-side prose) and
\texttt{ContractionSurjective} (the contraction is surjective;
equivalently, every linear functional on $\Lambda^2 E_-$ is the
contraction of $\omega$ with some vector --- \emph{primitive /
no spectator} in TFPT-side prose). Together with
\texttt{Nontrivial}\,$E_-$ they imply the Stage C contraction
iso and hence $\dim E_- = 3$.

\textbf{Where the standing Yukawa gap now lives.}
The Lean chain takes the trilinear form $\omega$ together with
the two predicates and \texttt{Nontrivial}\,$E_-$ as input.
The remaining open question is \emph{purely upstream}: does the
TFPT boundary kernel
(\texttt{CalderonCertificate} \emph{plus} representation-
theoretic data that the Yukawa coupling exists) produce a
trilinear form satisfying \texttt{ContractionInjective} and
\texttt{ContractionSurjective}? That question is TFPT theory
work, not Lean engineering. The Lean chain is now flat: there is
no hidden assumption between Stage D and $\dim E_- = 3$.

\subsection*{Layer 7 --- Mathlib bridge}

\begin{theorem}[\texttt{MathlibBridge.Polarization.toCompleteOrthogonalIdempotents}]
A \texttt{TFPT.Carrier.Polarization} on a ring $A$ is canonically
a Mathlib \texttt{CompleteOrthogonalIdempotents} on the pair
$(P_-, P_+)$. A converse map
\texttt{ofCompleteOrthogonalIdempotents} is also provided.
\end{theorem}

\subsection*{Layer 8 --- Bundled main theorems}

The carrier-rigidity payload of the project is the following set
of bundled theorems on
\texttt{CarrierPremises K M}:

\begin{theorem}[\texttt{CarrierPremises.primitive\_pair\_eq\_sm}]
Under the primitive oriented determinant-preserving carrier
condition with rank data $(m, n) = (3, 2)$, $(q_-, q_+) = (-2, 3)$.
\end{theorem}

\begin{theorem}[\texttt{CarrierPremises.trace\_Y\_eq\_zero}]
$\tr_M Y = 0$ for $Y = -\tfrac13 P_- + \tfrac12 P_+$.
\end{theorem}

\begin{theorem}[\texttt{CarrierPremises.carrier\_polynomial\_Y}]
$6\cdot(Y\cdot Y) - Y - \mathbf 1 = \mathbf 0$ in $\mathrm{End}_K M$.
The hypothesis $[\texttt{CharZero K}]$ is used here to provide
$\widehat 6 \in K$.
\end{theorem}

\begin{theorem}[\texttt{CarrierPremises.hypercharge\_carrier\_packet}]
Under \texttt{CarrierPremises} and the primitive oriented
determinant-preserving carrier condition, the three bundled
conclusions above hold \emph{simultaneously}. This is the
top-level claim of the project.
\end{theorem}

\subsection*{Concrete realisation}

\begin{theorem}[\texttt{Hypercharge.trace\_Y, Hypercharge.Y\_carrier\_polynomial}]
On the concrete $5\times 5$ rational diagonal model,
$Y = \mathrm{diag}(-\tfrac13,-\tfrac13,-\tfrac13,\tfrac12,\tfrac12)$
satisfies $\tr Y = 0$ and $6\,Y^2 - Y - \mathbf 1 = \mathbf 0$,
proved by direct per-entry rational arithmetic.
\end{theorem}

% ---------------------------------------------------------------- %
\section{Significance}
\label{sec:significance}
% ---------------------------------------------------------------- %

The two load-bearing chains, with all hypotheses explicitly
visible, are:

\paragraph{Chain~A --- trace identity.}
\[
\text{Calderón certificate}
\;\xRightarrow{\text{Paper 1 interface}}\;
\varepsilon^2 = \mathbf 1
\;\xRightarrow{\text{\texttt{InvolutionProjectors}}}\;
P_-, P_+ \text{ orthogonal idempotents}
\]
together with the rank discharges
\[
\dim E_+ = 2 \quad (\text{\texttt{HiggsIndexCertificate}}),
\qquad
\omega : \Lambda^3 E_- \to K\,\text{+CI+CS+Nontrivial} \quad (\text{Stage D})
\;\Rightarrow\;
\dim E_- = 3
\]
and the determinant-normalised coefficient choice
$Y := -\tfrac13 P_- + \tfrac12 P_+$ give, via
\texttt{TraceProjection},
\[
\tr Y = -\tfrac13 \cdot 3 + \tfrac12 \cdot 2 = 0.
\]
The Yukawa input is now an alternating trilinear form with two
explicit \texttt{Prop} predicates. The contraction iso and the
rank-3 conclusion are both derived theorems.

\paragraph{Chain~B --- integer rigidity.}
The trace-zero condition
\[
m\,q_- + n\,q_+ = 0
\]
follows from the universal-$\lambda$ determinant character. Combined
with the primitivity and orientation conditions
\[
q_- < 0 < q_+,\qquad
\gcd(|q_-|, q_+) = 1,
\]
the general $(m,n)$ rigidity theorem gives
\[
q_- = -\frac{n}{\gcd(m,n)},\qquad q_+ = \frac{m}{\gcd(m,n)}.
\]
For $(m,n) = (3,2)$ this specialises to $(q_-,q_+) = (-2,3)$.
The integer answer is a corollary of a parametric theorem; it is
not hard-coded.

\paragraph{Why both chains matter together.}
The carrier polynomial $6 Y^2 - Y - \mathbf 1 = \mathbf 0$ is a
consequence of the polarization layer alone (it does not need
either rank discharge). The trace identity $\tr Y = 0$ requires
the rank discharges. The integer pair $(-2, 3)$ requires the
primitive oriented determinant-preserving condition together
with $(m, n) = (3, 2)$. The bundled theorem
\texttt{hypercharge\_carrier\_packet} is the unique statement
that requires all upstream pieces simultaneously.

\paragraph{What is \emph{not} claimed.} The $3+2$ multiplet
$Y = \mathrm{diag}(-\tfrac13, -\tfrac13, -\tfrac13, \tfrac12, \tfrac12)$
is the carrier. The \emph{full Standard-Model hypercharge
spectrum} --- all chiral-multiplet hypercharges including duals,
exterior powers, and tensor products --- is not formalised in
this experiment. The title and main theorem both say "$3+2$
hypercharge carrier", not "SM hypercharge spectrum".

% ---------------------------------------------------------------- %
